% Fichier généré automatiquement par tools/generate_corrections.py.
% Source : STAT/STAT-02-Frottement/532_MAM_Frottement_Cylindre/532_MAM_Frottement_Cylindre.tex
% Statut : complété (4/4 question(s)).
\begin{corrigebox}[Corrigé complété]
\CorrigeQuestion{1}
Dans le plan $(\vec y,\vec z)$, on représente la génératrice du cylindre,
            la normale radiale $\vec e_r$ et la direction axiale $\vec z$. Dans le plan
            $(\vec x,\vec y)$, le point $Q$ est repéré par l'angle $\theta$ et
            $\vec e_r=\cos\theta\vec x+\sin\theta\vec y$. La pression de 2 sur 1 est
            dirigée vers l'axe, $-p\,dS\,\vec e_r$. À la limite d'un glissement de 1 vers
            $-\vec z$, l'effort d'adhérence de 2 sur 1 est dirigé vers $+\vec z$.
\CorrigeQuestion{2}
Avec $dS=R\,d\theta\,dz$, la loi de Coulomb à la limite d'adhérence donne
            $$\boxed{d\vec F_{2\to1}(Q)=-pR\,d\theta\,dz\,\vec e_r
            +fpR\,d\theta\,dz\,\vec z}.$$
            En régime d'adhérence stricte, la composante tangentielle a une norme
            inférieure ou égale à $fp\,dS$. L'intégration axiale sur toute la surface
            cylindrique conduit à $F_{z,\max}=2\pi f pRL$.
\CorrigeQuestion{3}
Exprimons le torseur des actions mécaniques sous sa forme locale en un point $M$ : 

$$
\torseurl{d\vectf{2}{1}}{d\vectm{M}{2}{1}=\vect{0}}{M}
$$

La forme globale au point O est alors donnée par :

$$
\torseurstat{T}{2}{1} = \torseurl{\vectf{2}{1} = \int d\vectf{2}{1}}{\vectm{M}{2}{1} = \int d\vectm{M}{2}{1}= \int \vect{OM}\wedge d\vectf{2}{1}}{M}
$$

\vspace{.5cm}

\textbf{Calculons $\vectf{2}{1}$.}

$$
\vectf{2}{1} = \int d\vectf{2}{1} = \iint p \vect{-r} dS = -p \iint  \vect{r} dS
= -p \iint  \left(\cos\theta\vect{x}+\sin\theta\vect{y} \right) dS $$

$$
\vectf{2}{1}
= -p \int\limits_{-L/2}^{L/2} \int\limits_{-\pi/2}^{\pi/2}   \left(\cos\theta\vect{x}+\sin\theta\vect{y} \right) Rd\theta dz
= -p R L \int\limits_{-\pi/2}^{\pi/2}   \left(\cos\theta\vect{x}+\sin\theta\vect{y} \right) d\theta 
$$
$$
\vectf{2}{1}
= -p R L \left(\int\limits_{-\pi/2}^{\pi/2}   \cos\theta\vect{x}d\theta + \int\limits_{-\pi/2}^{\pi/2} \sin\theta\vect{y} d\theta \right)
= -p R L \left(\left[\sin\theta \right]_{-\pi/2}^{\pi/2}\vect{x}
+\left[ -\cos\theta\right]_{-\pi/2}^{\pi/2}\vect{y}
\right)
$$

$$
\vectf{2}{1}
= -p R L \left(2\vect{x}
+0\vect{y}
\right) = -2pRL\vect{x}
$$

$2RL$ est appelée surface projetée du cylindre. Elle correspond au produit du diamètre par sa longueur.

\vspace{.5cm}

\textbf{Calculons $\vectm{M}{2}{1}$.}
$$
\vectm{M}{2}{1} = \int d\vectm{M}{2}{1}= \int \vect{OM}\wedge d\vectf{2}{1}
$$

$$
\vectm{M}{2}{1} = -p \iint R\vect{r} \wedge \vect{r}dS = \vect{0}
$$

Au final, 
$$
\torseurstat{T}{2}{1} = \torseurl{\vectf{2}{1} = -2pRL\vect{x}}{\vectm{M}{2}{1} =  \vect{0}}{M}
$$
\CorrigeQuestion{4}
Dans ce cas : 
$$
\vectf{2}{1} = \int d\vectf{2}{1} = \iint p(\theta) \vect{-r} dS 
= - p_0R\iint\cos\theta  \left(\cos\theta\vect{x}+\sin\theta\vect{y} \right)  d\theta dz$$

$$
\vectf{2}{1} 
= - p_0 L R\int\limits_{-\pi/2}^{\pi/2}\cos\theta  \left(\cos\theta\vect{x}+\sin\theta\vect{y} \right)  d\theta$$

$$
\int\limits_{-\pi/2}^{\pi/2}\cos^2\theta  d\theta = \dfrac{\pi}{2}
\quad 
\text{et}
\quad
\int\limits_{-\pi/2}^{\pi/2}\cos\theta \sin\theta  d\theta = 0
$$
Au final :
$$
\vectf{2}{1} 
= - p_0 L R \dfrac{\pi}{2}\vect{x}$$
\end{corrigebox}
